\documentclass[a4paper, oneside]{article}
\usepackage{graphicx}
\usepackage{amsthm}
\usepackage{amsmath}
\usepackage{amssymb}
\usepackage[a4paper,
            bindingoffset=0.2in,
            left=2cm,
            right=2cm,
            top=2cm,
            bottom=2cm,
            footskip=.25in]{geometry}
\usepackage[italian]{babel}
\usepackage{pgfplots}
\usepackage{tabularx}
\usepackage{tikz}
\usepackage{wrapfig}
\usepackage{color}
\usepackage[d]{esvect}
\definecolor{page}{rgb}{0.129,0.157,0.212}
\pagecolor{page}
\color{white}
\graphicspath{ {./images/} }
\usetikzlibrary{shapes.geometric}
\usetikzlibrary{datavisualization}
\usetikzlibrary{datavisualization.formats.functions}
\usetikzlibrary{patterns}
\pgfplotsset{width=10cm,compat=1.18}

\title{Appunti di Fisica II}
\author{Tommaso Miliani}
\date{16-03-26}

\begin{document}
\newtheoremstyle{theoremEnv}
                {}          % Space above
                {}          % Space below
                {\slshape}  % Body font
                {}          % Indent amount
                {\bfseries} % Head font
                {.}         % Punctuation after head
                {\newline}  % Space after theorem head
                {}          % Theorem head spec
\theoremstyle{theoremEnv}

\newtheorem{definition}{Definizione}[section]
\newtheorem{theorem}{Teorema}[section]
\newtheorem{lemma}{Proposizione}[section]
\newtheorem{observation}{Osservazione}[section]
\newtheorem{corollary}{Corollario}[theorem]
\newtheorem{example}{Esempio}[section]
\newtheorem{remark}{Enunciato}[section]

\maketitle

\section{Dielettrici}
Sono dei materiali che, se sottoposti ad un campo elettrico esterno, 
le cariche tendono a dipsorsi in modo tale che $\vv{p} = d\vv{E}$. Ossia 
si descrive il dielettrico come un materiale che ha una risposta lineare 
al campo elettrico a cui è sottoposto. In questo sistema non si è in grado di 
descrivere in maniera accurata il comportamento di un dielettrico ma si possono 
utilizzare delle relazioni empiriche per cui si individua un piccolo 
volume $d\tau$ per il dielettrico e si definisce la \textbf{densità di dipolo} 
per ogni elemento $d\tau$, in modo tale che
\begin{gather*}
    d\vv{p} = \vv{P}(\vv{r'}) d\tau  
\end{gather*}  
Il volumetto deve essere scelto piccolo a sufficienza in modo che la funzione $\vv{P}$ sia derivabile 
ma allo stesso tempo abbastanza grande in modo tale che lo possa definire come un insieme 
continuo di cariche. Il potenziale si ottiene con la formula dello sviluppo multipolare
\begin{gather*}
    dV(\vv{r}, \vv{r'}  ) =  \frac{1}{4\pi\epsilon_0} \frac{d\vv{p}(\vv{r} - \vv{r'}  )}{\left| \vv{r} - \vv{r'}   \right|^{3} } 
\end{gather*}
Adesso si pul sostituire $d\vv{p}$:
\begin{gather*}
    dV(\vv{r}, \vv{r'}  ) = \frac{1}{4\pi \epsilon_0} \frac{\vv{P}(\vv{r} ) d\tau (\vv{r} - \vv{r'}  )}{\left| \vv{r} - \vv{r'}   \right|^{3} }
\end{gather*} 
Integrando ora su tutto il volume (mantenendo fisso $\vv{r}$ ) si ottiene che
\begin{gather*}
    V(\vv{r}, \vv{r'}  ) = \frac{1}{4\pi\epsilon_0} \int_{\tau}  \frac{\vv{P}(\vv{r} ) (\vv{r} - \vv{r'}  ) }{\left| \vv{r} - \vv{r'}   \right|^{3} } d\tau
\end{gather*}
Si è già visto che
\begin{gather*}
    \frac{\vv{r} - \vv{r'}  }{\left| \vv{r} - \vv{r'}   \right|^{3} } = \vv{\nabla} \frac{1}{\left| \vv{r} - \vv{r'}   \right| } \qquad R = \left| \vv{r} - \vv{r'}   \right|   
\end{gather*}
Dunque
\begin{gather*}
    V(\vv{r} ,\vv{r'}  ) = \frac{1}{4\pi\epsilon_0} \int_{\tau} \vv{P} \cdot \vv{\nabla} \frac{1}{R} = \frac{1}{4\pi\epsilon_0} \int_{\tau}\left(\vv{\nabla}\left(\frac{\vv{P} }{R}\right)  - \frac{1}{R} \vv{\nabla} \cdot \vv{P}  \right)  d\tau
\end{gather*}
Dove
\begin{gather*}
    \sum_{i = 1}^{3} P_i \frac{\partial }{\partial x_i}  \frac{1}{R} =    \sum_{i = 1}^{3} \frac{\partial }{\partial x_i} \left(\frac{P_i}{R}\right)  = \vv{\nabla}\left(\frac{\vv{P} }{R}\right)  - \frac{1}{R} \vv{\nabla} \cdot \vv{P}  
\end{gather*}
Allora, ricordando lo sviluppo al primo ordine del multipolare 
\begin{gather*}
    V(\vv{r} ) = \frac{1}{4\pi\epsilon_0} \int_{\tau} \frac{\rho(\vv{r'} )}{\left| \vv{r} - \vv{r'}   \right| } d\tau + \frac{1}{4\pi\epsilon_0} \int_{S} \frac{\sigma(\vv{r'} )}{\left| \vv{r} - \vv{r'}   \right| } dS  
\end{gather*}
Si può ora esprimere un campo scalare, dentro il dielettrico, che prende il nome di - ((MANCA PEZZO))


Se si mette un cubo dielettrico dentro ad un condensatore 
il campo elettrico diminuisce poiché le cariche si dispongono 
opposte risoetto alle cariche sulle piastre del condensatore. 
(Manca un pezzo)
Dalla prima di Maxwell
\begin{gather*}
    \vv{\nabla} \cdot \vv{E} = \frac{\rho}{\epsilon_0} = \frac{\rho_L + \rho_P}{\epsilon_0} = \frac{\rho_L}{\epsilon_0} - \frac{1}{\epsilon_0} \vv{\nabla} \cdot \vv{P}    
\end{gather*}
SI definisce il vettore spostamento come $\vv{D} = \epsilon_0 \vv{E} + \vv{p}  $ e 
allora 
\begin{align}
    \vv{\nabla} \cdot \vv{D} = \rho_L  
\end{align}   
Ci si limita a studiare materiali che sono 
\begin{itemize}
    \item Lineari: la risposta al campo elettrico è lineare
    \item Isotropi: la risposta al campo non cambia se cambia la direzione del campo
    \item Omogenei: tutti i punti del materiale sono uguali
    \item Locali: la risposta non dipende da punto a punto
\end{itemize}

\end{document}